\documentclass[aspectratio=169]{beamer}

% -------------------- Packages --------------------
\usepackage{graphicx}
\usepackage{amsmath}
\usepackage{amssymb}
\usepackage{booktabs}
\usepackage{times}
\usepackage{listings}
\usepackage{xcolor}
\usepackage{tikz}

% -------------------- Theme --------------------
\usetheme{Boadilla}
\usecolortheme{rose}

% -------------------- Custom Colors --------------------
\definecolor{securityblue}{RGB}{30, 55, 110}
\definecolor{accentblue}{RGB}{52, 152, 219}
\definecolor{darkred}{RGB}{180, 30, 30}

\setbeamercolor{structure}{fg=securityblue}
\setbeamercolor{palette primary}{bg=securityblue, fg=white}
\setbeamercolor{palette secondary}{bg=accentblue, fg=white}
\setbeamercolor{palette tertiary}{bg=securityblue, fg=white}
\setbeamercolor{frametitle}{bg=securityblue, fg=white}
\setbeamercolor{alerted text}{fg=darkred}

\setbeamertemplate{caption}[numbered]

% -------------------- Footer --------------------
\setbeamertemplate{footline}{
  \leavevmode
  \hbox{
  \begin{beamercolorbox}[wd=\paperwidth, ht=3ex, dp=1.2ex, leftskip=0.4cm, rightskip=0.4cm]{palette primary}
  \color{white}\scriptsize
  SPi: AI-Based Insider Threat Detection System \;|\; AJIN JOSE - Roll No: 07 \hfill \insertframenumber/17
  \end{beamercolorbox}}
}

% -------------------- Title Page Setup --------------------
\title{SPi: AI-Based Insider Threat Detection System}
\author{AJIN JOSE \\ Roll No: 07 \\ Reg No: KTE24MCA--2008}
\institute{Rajiv Gandhi Institute of Technology (RIT), Kottayam \\ Department of Computer Applications}
\date{\today}

\begin{document}

% -------------------- TITLE SLIDE --------------------

\begin{frame}[plain]
\centering
\vspace{0.5cm}
\begin{beamercolorbox}[sep=12pt,center,shadow=true,rounded=true]{palette primary}
\fontsize{18pt}{20pt}\selectfont 
\textbf{SPi: AI-Based Insider Threat Detection System}
\end{beamercolorbox}

\vspace{1.5cm}
{\Large \textbf{AJIN JOSE}}\\[0.3cm]
{\normalsize Roll No: 07}\\
{\normalsize Reg No: KTE24MCA--2008}

\vspace{1cm}
{\small Guided By: Prof. Jincy Kuriakose}\\
{\small Department of Computer Applications}\\
{\small Rajiv Gandhi Institute of Technology (RIT), Kottayam}

\vspace{1.5cm}
{\tiny February 2026}
\end{frame}

% -------------------- TABLE OF CONTENTS --------------------

\begin{frame}{Table of Contents}
\tableofcontents
\end{frame}

% -------------------- SECTION: INTRODUCTION --------------------

\section{Introduction \& Objectives}

\begin{frame}{1. Introduction}

\begin{itemize}
    \item \textbf{Insider Threats:} Malicious or negligent actions by authorized users
    \item \textbf{Challenge:} Traditional security systems focus on external attacks
    \item \textbf{Solution:} Behavioral anomaly detection and risk scoring
    \vspace{0.3cm}
    \item \textbf{SPi Integration:}
    \begin{itemize}
        \item Digital activity logs (file operations, logins, USB access)
        \item Physical access records (CCTV, entry/exit logs)
        \item Machine learning–based anomaly detection (Isolation Forest)
        \item Dynamic quantitative Risk Score
    \end{itemize}
\end{itemize}

\end{frame}

\begin{frame}{2. Objectives}

\begin{itemize}
    \item[\checkmark] Identify malicious or negligent activities through behavioral log analysis
    \vspace{0.2cm}
    \item[\checkmark] Correlate digital system events with physical access through computer vision
    \vspace{0.2cm}
    \item[\checkmark] Apply \textbf{Isolation Forest} techniques to detect deviations from normal behavior
    \vspace{0.2cm}
    \item[\checkmark] Support proactive security decisions via \textbf{dynamic risk scoring}
    \vspace{0.2cm}
    \item[\checkmark] Provide real-time threat monitoring and alerting dashboard
    \vspace{0.2cm}
    \item[\checkmark] Integrate behavioral and CCTV-based detection for comprehensive threat assessment
\end{itemize}

\end{frame}

% -------------------- SECTION: SYSTEM DESIGN --------------------

\section{System Design \& Architecture}

\begin{frame}{3. System Architecture}

\begin{figure}
    \centering
    \includegraphics[height=0.65\textheight]{Untitled Diagram.png} 
    \caption{Data Flow: Preprocessing to Dashboard}
\end{figure}

\end{frame}

\begin{frame}{3.1 System Components}

\begin{itemize}
    \item \textbf{Data Ingestion Layer}
    \begin{itemize}
        \item Employee behavioral logs (CSV uploads)
        \item CCTV video processing
        \item Face recognition integration
    \end{itemize}
    \vspace{0.3cm}
    \item \textbf{Anomaly Detection Engine}
    \begin{itemize}
        \item Isolation Forest algorithm
        \item Feature normalization
        \item Risk score calculation
    \end{itemize}
    \vspace{0.3cm}
    \item \textbf{Analytics \& Dashboard}
    \begin{itemize}
        \item Risk visualization (scatter plots, distributions)
        \item Real-time monitoring
        \item Alert management
    \end{itemize}
\end{itemize}

\end{frame}

% -------------------- SECTION: MACHINE LEARNING --------------------

\section{Machine Learning Model}

\begin{frame}{4. Isolation Forest – Core Concept}

\begin{columns}[T]

\begin{column}{0.5\textwidth}
\begin{itemize}
    \item \textbf{Unsupervised} anomaly detection
    \item Random partitioning isolates anomalies faster than normal points
    \item Anomalies require fewer recursive splits
    \item Builds multiple random decision trees
\end{itemize}
\end{column}

\begin{column}{0.5\textwidth}
\[
\text{Path Length} \propto \text{Anomaly Score}
\]

\vspace{0.5cm}
\begin{itemize}
    \item Short path → High anomaly probability
    \item Long path → Normal behavior
\end{itemize}
\end{column}

\end{columns}

\end{frame}

\begin{frame}{5. Isolation Forest – Mathematical Foundation}

\textbf{Expected Path Length (Normalization Factor):}

\[
c(n) = 2H(n-1) - \frac{2(n-1)}{n}
\]

{\small Where: $H(n)$ = Harmonic number, $n$ = Total samples}

\vspace{0.5cm}

\textbf{Anomaly Score:}

\[
s(x,n) = 2^{-\frac{E(h(x))}{c(n)}}
\]

{\small Where: $E(h(x))$ = Average path length of sample $x$}

\vspace{0.5cm}

\begin{itemize}
    \item Score range: $[0, 1]$
    \item $s \approx 1$ → Strong anomaly
    \item $s \approx 0.5$ → Normal behavior
    \item $s \approx 0$ → Deep normal instance
\end{itemize}

\end{frame}

\begin{frame}{6. Isolation Forest Algorithm}

\begin{enumerate}
    \item \textbf{Input:} Dataset $D$ with $n$ samples
    \vspace{0.2cm}
    
    \item \textbf{For each tree:}
    \begin{itemize}
        \item Randomly sample subset of data
        \item Recursively partition using random feature and split value
        \item Build isolation tree until all points isolated
    \end{itemize}
    \vspace{0.2cm}
    
    \item \textbf{Compute:} Average path length for each sample across all trees
    \vspace{0.2cm}
    
    \item \textbf{Generate:} Anomaly score using normalization formula
    \vspace{0.2cm}
    
    \item \textbf{Output:} Anomaly scores for threshold-based classification
\end{enumerate}

\end{frame}

% -------------------- SECTION: RISK MODEL --------------------

\section{Risk Scoring Model}

\begin{frame}{7. Risk Score Mathematical Model}

\textbf{Integrated Risk Formula:}

\[
R = w_1 A_d + w_2 A_p + w_3 B_f
\]

\vspace{0.3cm}

\textbf{Component Definitions:}
\begin{itemize}
    \item $A_d$ = Digital anomaly score (Isolation Forest)
    \item $A_p$ = Physical access risk (CCTV-based detection)
    \item $B_f$ = Behavioral frequency factor (anomaly event density)
\end{itemize}

\vspace{0.3cm}

\textbf{Constraint:}
\[
w_1 + w_2 + w_3 = 1
\]

{\small Default: $w_1 = 0.5$, $w_2 = 0.3$, $w_3 = 0.2$}

\end{frame}

\begin{frame}{8. Digital Anomaly Score}

\textbf{Feature Components:}
\begin{itemize}
    \item Login patterns and frequency
    \item Night-time login activity
    \item File operations (accessed, modified, deleted, copied)
    \item USB device connections
    \item External email communications
    \item Database access patterns
\end{itemize}

\vspace{0.3cm}

\textbf{Normalization:}
\[
A_d = \frac{\text{ISO Forest Score} - \text{Min}}{\text{Max} - \text{Min}}
\]

Produces normalized risk scale $[0, 100]$

\end{frame}

\begin{frame}{9. Physical Access Risk}

\textbf{CCTV-Based Risk Assessment:}

\[
A_p = \frac{\text{Unauthorized Access Events}}{\text{Total Access Attempts}}
\]

\vspace{0.3cm}

\textbf{Factors Considered:}
\begin{itemize}
    \item Unauthorized zone entries
    \item Detection confidence levels
    \item Extended presence duration
    \item After-hours facility access
    \item Restricted area movements
\end{itemize}

\vspace{0.3cm}

Higher ratio → Increased suspicious physical movement

\end{frame}

\begin{frame}{10. Behavioral Frequency Factor}

\textbf{Definition:}

\[
B_f = \frac{\text{Anomalous Events Detected}}{\text{Total Events}}
\]

\vspace{0.3cm}

\textbf{Examples of Anomalous Events:}
\begin{itemize}
    \item Multiple after-hours logins
    \item Bulk file downloads/deletions
    \item USB mass storage operations
    \item External data transfers
    \item Unusual database queries
    \item Night-time system access patterns
\end{itemize}

\vspace{0.3cm}

Captures frequency and persistence of risky activities

\end{frame}

\begin{frame}{11. Risk Level Classification}

\begin{table}
\centering
\renewcommand{\arraystretch}{1.5}
\begin{tabular}{|c|c|l|}
\hline
\textbf{Risk Score Range} & \textbf{Level} & \textbf{Action} \\
\hline
0.00 – 0.30 & \textcolor{green!70!black}{\textbf{Low}} & Monitor routine \\
\hline
0.31 – 0.60 & \textcolor{orange}{\textbf{Medium}} & Increase monitoring \\
\hline
0.61 – 0.80 & \textcolor{blue}{\textbf{High}} & Flag for review \\
\hline
0.81 – 1.00 & \textcolor{red}{\textbf{Critical}} & Immediate investigation \\
\hline
\end{tabular}
\end{table}

\vspace{0.5cm}
System generates \textbf{real-time alerts} when thresholds are exceeded.

\end{frame}

% -------------------- SECTION: IMPLEMENTATION --------------------

\section{Implementation \& Results}

\begin{frame}{12. Implementation Details}

\textbf{Technology Stack:}
\begin{itemize}
    \item \textbf{Frontend:} React.js with TypeScript, Tailwind CSS
    \item \textbf{Backend:} Python Flask/FastAPI
    \item \textbf{ML Engine:} Scikit-learn (Isolation Forest)
    \item \textbf{CCTV:} CompreFace API for face recognition
    \item \textbf{Database:} Local storage with CSV imports
\end{itemize}

\vspace{0.3cm}

\textbf{Model Configuration:}
\begin{itemize}
    \item Isolation Forest: 300 trees
    \item Contamination parameter: 0.05
    \item Feature normalization: StandardScaler
    \item Dataset: CERT Insider Threat (modified)
\end{itemize}

\end{frame}

\begin{frame}{13. Results: Threat Monitoring Dashboard}

\setlength{\columnsep}{12pt}

\begin{columns}[T]

\begin{column}{0.33\textwidth}
\begin{figure}
\centering
\includegraphics[width=\linewidth]{Screenshot 2026-02-15 203629.png}
\caption{Risk Management Dashboard}
\end{figure}
\end{column}

\begin{column}{0.33\textwidth}
\begin{figure}
\centering
\includegraphics[width=\linewidth]{Screenshot 2026-02-15 203620.png}
\caption{Individual Risk Profile}
\end{figure}
\end{column}

\begin{column}{0.33\textwidth}
\begin{figure}
\centering
\includegraphics[width=\linewidth]{Screenshot 2026-02-15 203608.png}
\caption{Scatter Plot Analysis}
\end{figure}
\end{column}

\end{columns}

\end{frame}

\begin{frame}{14. Key Features Implemented}

\begin{itemize}
    \item [\checkmark] \textbf{Data Ingestion:} CSV upload for behavioral data
    \vspace{0.2cm}
    
    \item [\checkmark] \textbf{CCTV Integration:} Face detection and unauthorized access flagging
    \vspace{0.2cm}
    
    \item [\checkmark] \textbf{Risk Assessment:} Real-time scoring with multi-component analysis
    \vspace{0.2cm}
    
    \item [\checkmark] \textbf{Analytics Dashboard:} Interactive visualizations and correlations
    \vspace{0.2cm}
    
    \item [\checkmark] \textbf{Alert System:} Email/SMS notifications for critical threats
    \vspace{0.2cm}
    
    \item [\checkmark] \textbf{Authorized Face Database:} Special emp1/emp2 mapping for high-risk scenarios
\end{itemize}

\end{frame}

% -------------------- SECTION: PROJECT PLAN --------------------

\section{Project Timeline}

\begin{frame}{15. Project Plan \& Development Timeline}

\small
\begin{table}
\centering
\renewcommand{\arraystretch}{1.4}
\begin{tabular}{|p{2cm}|p{5.5cm}|p{2cm}|}
\hline
\textbf{Phase} & \textbf{Work Items} & \textbf{Status} \\
\hline
\textbf{Phase 1} & Literature review, CERT dataset analysis, threat modeling & Completed \\
\hline
\textbf{Phase 2} & System architecture, UI design, backend API development & Completed \\
\hline
\textbf{Phase 3} & Isolation Forest integration, risk scoring implementation & Completed \\
\hline
\textbf{Phase 4} & Face recognition, CCTV correlation, dual-source detection & Completed \\
\hline
\textbf{Phase 5} & Testing, documentation, deployment preparation & Pending \\
\hline
\end{tabular}
\caption{Development Timeline Summary}
\end{table}

\end{frame}

% -------------------- SECTION: CONCLUSION --------------------

\section{Conclusion \& Future Work}

\begin{frame}{16. Conclusion}

\textbf{Achievements:}
\begin{itemize}
    \item Integrated digital and physical anomaly detection
    \item Designed comprehensive mathematical risk scoring framework
    \item Implemented scalable system for enterprise deployment
    \item Real-time monitoring with multi-component risk assessment
\end{itemize}

\vspace{0.4cm}

\textbf{Key Contributions:}
\begin{itemize}
    \item Novel correlation between behavioral logs and CCTV analysis
    \item Weighted risk model addressing multiple threat vectors
    \item User-friendly dashboard for security analysts
\end{itemize}

\vspace{0.4cm}

\textbf{Future Enhancements:}
\begin{itemize}
    \item Time-series modeling for temporal anomalies
    \item Deep learning for pattern recognition
    \item Graph-based relationship analysis
    \item Multi-organization threat correlation
\end{itemize}

\end{frame}

\begin{frame}{17. Development History}

\begin{figure}
    \centering
    \includegraphics[width=0.7\textwidth]{Screenshot 2026-02-15 201548.png}
    \caption{Git Repository Development History and Commits}
\end{figure}

\end{frame}

% -------------------- BIBLIOGRAPHY --------------------

\section{References}

\begin{frame}{18. Bibliography \& Research References}

\begin{thebibliography}{99}

\bibitem{Liu2008}
F. T. Liu, K. M. Ting, and Z. H. Zhou, 
``Isolation Forest,'' 
\textit{IEEE International Conference on Data Mining (ICDM)}, 
pp. 413-422, 2008.

\bibitem{Insider2016}
M. M. Tuor, S. Kaplan, B. Moskovitch, C. Shahar, and A. Feiglin,
``Challenging the Insider Threat Detection Paradigm Using Extreme Value Theory,''
\textit{Journal of Cybersecurity}, 
vol. 2, no. 1, pp. 39-62, 2016.

\bibitem{Chandola2009}
V. Chandola, A. Banerjee, and V. Kumar,
``Anomaly Detection: A Survey,''
\textit{ACM Computing Surveys}, 
vol. 41, no. 3, pp. 1-58, 2009.

\bibitem{Maxion2010}
R. A. Maxion and R. E. Rehfuss,
``Improving User Masquerade Detection and Identification with Hybrid-KNN and Behavioral Knowledge,''
\textit{International Journal of User Authentication and Authentication}, 
pp. 123-145, 2010.

\bibitem{Kamran2019}
K. R. Patel and B. M. Mehta,
``Insider Threat Detection Using Machine Learning Techniques: A Review,''
\textit{International Journal of Advanced Trends in Computer Science and Engineering}, 
vol. 8, no. 3, pp. 567-579, 2019.

\end{thebibliography}

\end{frame}

% -------------------- THANK YOU SLIDE --------------------

\begin{frame}[plain]
\centering
\vspace{3cm}
{\Huge \textbf{Thank You}}

\vspace{1.5cm}

{\Large Questions?}

\vspace{1cm}

{\small AJIN JOSE | Roll No: 07 \\ Rajiv Gandhi Institute of Technology, Kottayam}
\end{frame}

\end{document}
